
\exo

\begin{enumerate}


\item Dans une classe de seconde de 30 élèves, on sait que 6 d'entre-eux font du roller.
\

Calculer la proportion d'élèves de cette classe pratiquant le roller.
\

Exprimer le résultat sous la forme d'une fraction irréductible et interpréter le résultat à l'aide d'une phrase.


\item  Dans une classe de Première de $32$ élèves, le pourcentage d'élève pratiquant le roller est de $31,25 \%$.
\

Calculer le nombre d'élèves de cette classe pratiquant le roller.

\item Les 128 élèves du lycée qui pratiquent le roller représentent $22 \%$ du total des élèves du lycée.
\

Calculer le nombre total d'élèves du lycée.

\item Dans un camping à Porto-Vecchio, il y a $380$ français, $170$ anglais, $400$ hollandais et $30$ belges.

Quelle est la proportion de chaque nationalité dans ce camping ? 

Exprimer le résultat sous la forme d'une fraction irréductible, puis en pourcentage (arrondi au dixième).

\item En 2$016$, la France compte environ $67$ millions d'habitants. Arrondir à $0,01$ millions près.

a) $20\%$ des Français ne partent pas en vacances. Combien de Français ne partent-ils pas en vacances ?

b) Les deux tiers des vacanciers français partent à l'étranger. Combien de Français partent-ils à l'étranger ?

\item a) Rose gagne $2200 \EUR$ par mois. Le montant de son loyer équivaut à $40\%$ de son salaire et le crédit pour sa voiture équivaut à 25\% de son salaire. Quel est le montant du loyer et du remboursement du crédit ?

b) Rose vit en couple avec son mari. Le salaire de Rose représente $60\%$ des revenus du couple.

Calculer alors le montant total des revenus du couple et en déduire le salaire du mari de Rose.


\end{enumerate}

\exo

 Dans une classe, $45\%$ des élèves sont des garçons. Parmi eux, $60\%$ pratiquent une activité sportive. 

Quelle est la proportion de garçons pratiquant du sport dans l'ensemble des élèves de la classe ?

\medskip

\exo 

Une enquête auprès d'élèves fumeurs montre que 90\% d'entre eux ont déjà essayé d'arrêter de fumer et que parmi ces dernier, 60\% ont réussi à s'arrêter plus d'un mois. Calculer la proportion d'élèves ayant réussi à s'arrêter plus d'un mois parmi les élèves fumeurs interrogés.

\medskip

\exo 

Dans un lycée, 92\% des élèves de seconde sont passés en classe de première.

Parmi ces élèves la proportion d'élèves qui sont passés en première STMG est égale à 0,35.

Calculer la proportion d'élèves qui sont passés en première TMG parmi les élèves qui étaient en seconde.

\medskip


\exo 

Une étude menée sur les groupes sanguins dans une population a établi que 42\% des individus de cette population sont du groupe O et que, en particulier, 36\% sont du groupe O avec le rhésus + (noté O+).

Calculer la fréquence des individus du groupe O+ dans la population des individus du groupe O.

\medskip

\exo

Au premier janvier 2011, il y avait en France métropolitaine 16269 personnes âgées de 100 ans et plus.

Parmi elles, 88,85\% étaient des femmes et 6232 femmes avaient 100 ans.

Calculer la proportion des femmes âgées de 100 ans parmi les femmes âgées de 100 et plus.

\medskip

\exo 

Une enquête auprès d'un groupe de personnes a établi que parmi elles :

- 82\% ont au moins 18 ans ;

- 42\% sont titulaires du permis de conduire automobile.

Calculer la proportion de personnes titulaires du permis de conduire automobile parmi les personnes du groupe ayant au moins 18 ans.


\exo 

 En  France,  le  prix  moyen  de  la  baguette  est  passé  de 0.70 \euro en 2005 à 0.92 \euro en 2015.

Calculer  la  variation  absolue  et  le  taux  d'évolution  du  prix  moyen  de  la  baguette  entre 2005 et 2015.

\medskip

\exo 

Une famille a consommé 150 mètres cubes d'eau en 2017 et 137 mètres cubes en 2018.

Calculer la variation relative (taux d'évolution) de la consommation de cette famille en 2017 à sa consommation en 2018.

\exo 

Un théâtre a programmé 260 représentations cette année contre 240 l'an dernier.

Calculer le taux d'évolution du nombre de représentations de l'année dernière à cette année.

\exo 

Au siècle dernier la population mondiale a triplé en 80 ans passant à 6 milliards d'individus.

Déterminer le taux d'évolution du nombre de Terriens sur ces 80 ans.

\exo

\begin{enumerate}
\item Donner les Coefficients multiplicateurs associés aux augmentations de : 

14\% \qquad\qquad 5.5\%\qquad\qquad 30\% \qquad\qquad 300\% \qquad\qquad 25,7\%

\item Donner les Coefficients multiplicateurs associés aux diminutions de : 

4\% \qquad\qquad 15\%\qquad\qquad 25.8\% \qquad\qquad 78.9\% \qquad \qquad 19.6\%
\end{enumerate}

\exo

 Le nombre de licenciés dans un club de sport augmente régulièrement de 4\% par an.

En 2015, il y avait 2540 licenciés.

\begin{enumerate}
\item Quel est le nombre de licenciés en 2016 ? en 2017 ?
\item Quel était le nombre de licenciés en 2014 ? et en 2013 ?
\end{enumerate}




\exo

Un magasin solde tous ces articles à 15\%.

\begin{enumerate}
\item Quel est le prix soldé d'un article qui valait 132\euro.
\item Calculer le prix initial d'un article soldé 35,70\euro .
\end{enumerate}

\exo 

\begin{enumerate}
\item  À quelle évolution globale correspond une augmentation de $10\%$ suivie d'une augmentation de $35\%$ ?
\item  À quelle évolution globale correspond une diminution de $50\%$ suivie d'une diminution de $20\%$ ?

\item  À quelle évolution globale correspond une augmentation de $5\%$ suivie d'une diminution de $18\%$ ?

\item À quelle évolution globale correspondent trois augmentations successives de $20\%$ ?

\item À quelle évolution globale correspondent quatre diminutions successives de $5 \%$ ?
\end{enumerate}

\exo 

 Le prix d'une action initialement  à $21,17\EUR$ subit une hausse de $3.4\%$ puis une diminution de $5.2\%$.

Calculer le nouveau prix de l'action au terme de ces deux évolutions. 

\exo 

 Après avoir subi une première baisse de $40\%$ puis une seconde de $10\%$, le prix d'un article s'élève à $217\EUR$.

Quelle était l'ancien prix de cet article avant les deux diminutions ?

\exo 

 Après une hausse de $12\%$ et une évolution inconnue, le prix d'un article passe de $47\EUR$ à $50\EUR$.

Quelle est la valeur de l'évolution mystère ?

\exo 

 Une valeur qui augmente de $30\%$ puis diminue de $30\%$ revient-elle à sa valeur initiale ? 

\exo 

Donner les taux réciproques de :

\begin{multicols}{3}
\begin{enumerate}
\item une hausse de 20 \% 
\item une baisse de 40 \%
\item une hausse de 5,5 \%
\item une baisse de 19,7 \%
\item une hausse de 38 \%
\item une baisse de 0,5 \%
\end{enumerate}
\end{multicols}

\exo

Le chiffre d'affaires annuel d'un hypermarché était en 2016 de $\np{35680200}$\euro ~ et, en 2017, de $\np{24872400} $ \euro ~.

\begin{enumerate}

\item Calculer le taux d'évolution, exprimé en pourcentage, du chiffre d'affaires entre 2016 et 2017. 
\item Calculer le pourcentage de la hausse qui ramènerait, en 2018, le chiffre d'affaires au niveau de celui de 2016.

\end{enumerate}

\exo 

\emph{Cet exercice est un questionnaire à choix multiple (QCM).\\ 
Pour chaque question, plusieurs réponses sont proposées, parmi lesquelles une seule est correcte et aucune  justification n'est demandée.\\ 
Chaque bonne réponse rapporte un point (sur 20), une question sans réponse ou fausse ne rapporte aucun point.}

\smallskip

\begin{alterqcm}[lq=10.5cm]
\AQquestion{Lors des soldes, une paire de chaussures porte l'étiquette  suivante :\\
 "Première démarque : $-20$\,\% puis démarque supplémentaire : $-10$\,\%"\\
Le taux d'évolution global associé au prix de la paire est une baisse de  : }{ 28\% ,  11\text{,}8\% ,  30\%}
\AQquestion{La valeur d'une imprimante achetée 850~\euro{} se déprécie de 20\,\% par an.\\ 
Quelle est sa valeur après trois ans ?}{340~\euro , 435\text{,}20~\euro , 544~\euro , 498\text{,}85~\euro}
\AQquestion{ Les dépenses du service Communication d'une entreprise sont passées de \np{2000}~\euro{} en 2005 à \np{6800}~\euro{} en 2008. \\
Le pourcentage d'augmentation est :}{3\text{,}4\% , 340 \% , 240 \% , 48 \%}

\end{alterqcm}

\smallskip
\exo 
 
Le tableau suivant donne le taux d'inflation annuel des prix en Chine entre 2007 et 2014: 


\begin{center}
\renewcommand{\arraystretch}{1.35}
\setlength{\tabcolsep}{0.28cm}
\begin{tabular}{|>{\bfseries}m{3.7cm}|*{7}{>{\centering\arraybackslash}m{1.08cm}|}}
\hline
Année & 2007 & 2008 & 2009 & 2010 & 2011 & 2012 & 2013 \\
\hline
Taux d'inflation en \% & 6,6 & 5 & $-1,2$ & 3,2 & 6,4 & 2 & 2,5 \\
\hline
\end{tabular}
\end{center}

\smallskip 
 
On admet que chaque année le taux d'évolution de la valeur de cette marchandise est égal au taux d'inflation en Chine.
 
Par exemple  le taux d'évolution de la valeur d'une  marchandise entre le 01/01/2007 et le 01/01/2008 
était $6,6$\,\%.
	\begin{enumerate}
		\item Une marchandise produite en Chine valait  \np{1500}~euros au 01/01/2007. Calculer la valeur de cette marchandise le 01/01/2008. 
		\item Une autre marchandise valait \np{2000}~euros au 01/01/2014, Calculer sa valeur au 01/01/2013.
		\item Quel est le taux d'évolution entre le 01/01/2007 et le 01/01/2014 ?
		\item Calculer, en pourcentage, à 0,1\,\% près, le taux d'évolution global de la valeur de la marchandise au cours des deux années comprises entre le 01/01/2010 et le 01/01/2012.
		\item Déterminer le taux annuel moyen entre le 01/01/2010 et le   01/01/2012 c'est à dire le taux $t\%$ d'augmentation qu'il faut prendre les deux années afin d'avoir une évolution équivalente à celle donnée dans le tableau sur la même période.
	\end{enumerate} 
